\documentclass{amsart}

\usepackage{amssymb}
\usepackage{booktabs}
\usepackage[hyphens]{url}
\usepackage[
  pdftitle={Embedding dimension strata in a space of complete subalgebras of power series},
  pdfauthor={Knowledge Sense},
  pdfsubject={Mathematics Subject Classification 14H20, 13J05, 20M14},
  pdfkeywords={complete subalgebra, formal power series, numerical semigroup, curve singularity, embedding dimension, normal form},
  colorlinks=true, linkcolor=black, citecolor=black, urlcolor=black
]{hyperref}

\newtheorem{theorem}{Theorem}[section]
\newtheorem{proposition}[theorem]{Proposition}
\newtheorem{lemma}[theorem]{Lemma}
\newtheorem{corollary}[theorem]{Corollary}
\theoremstyle{definition}
\newtheorem{definition}[theorem]{Definition}
\newtheorem{example}[theorem]{Example}
\newtheorem{question}[theorem]{Question}
\theoremstyle{remark}
\newtheorem{remark}[theorem]{Remark}

\numberwithin{equation}{section}

\newcommand{\CC}{\mathbb{C}}
\newcommand{\NN}{\mathbb{N}}
\newcommand{\ZZ}{\mathbb{Z}}
\newcommand{\AF}{\mathbb{A}}
\newcommand{\Rs}{\mathcal{R}_\Gamma}
\newcommand{\Rd}[1]{\mathcal{R}^{(#1)}_\Gamma}
\newcommand{\Ed}[1]{\mathcal{E}_{#1}}
\newcommand{\m}{\mathfrak{m}}
\newcommand{\ps}{[\![t]\!]}
\DeclareMathOperator{\ord}{ord}
\DeclareMathOperator{\edim}{edim}
\DeclareMathOperator{\rank}{rank}
\DeclareMathOperator{\wt}{wt}
\DeclareMathOperator{\red}{red}
\DeclareMathOperator{\Aut}{Aut}
\DeclareMathOperator{\Spec}{Spec}
\DeclareMathOperator{\spn}{span}

\begin{document}

\title[Embedding dimension strata]{Embedding dimension strata in a space of complete subalgebras of power series}

\author{Knowledge Sense}
\address{Knowledge Sense}
\email{math@knowledgesense.jp}

\date{12 October 2026}

\subjclass[2020]{Primary 14H20; Secondary 13J05, 20M14}
\keywords{Complete subalgebra, formal power series, numerical semigroup, curve singularity, embedding dimension, normal form}

\begin{abstract}
Hamilton parametrized the complete $\CC$-subalgebras of $\CC\ps$ with a prescribed value semigroup $\Gamma$ by an affine variety $\Rs$. Hamilton asked whether the locus $\Rd{d}$ of subalgebras generated by at most $d$ elements is cut out by the nonvanishing of leading coefficients of fixed polynomials, and whether it has the form $(\CC^*)^{g+1-d}\times\CC^{N-g+d-1}$ when $\Rs\cong\CC^N$. For $\Gamma=\langle 10,11,34,46,58\rangle$ we determine $\Rs$ and all its embedding dimension strata: there are global polynomial coordinates $(q,A,B,C)$ identifying $\Rs$ with $\AF^{54}\times\AF^3$, in which the embedding dimension of a subalgebra equals $5-\rank\left(\begin{smallmatrix}A&C\\0&B\end{smallmatrix}\right)$. Consequently $\Rd{4}\cong\AF^{54}\times(\AF^3\setminus\{0\})$. This smooth open subset of $\Rs\cong\AF^{57}$ is not affine, is not the nonvanishing locus of a single polynomial, and is not even homeomorphic to $\CC^*\times\CC^{56}$. Hence no fixed polynomial detects $\Rd{4}$ through the leading coefficient of its reduction, as asked in Hamilton's Question~4.7, and $\Rd{4}$ is not isomorphic to the space $\CC^*\times\CC^{56}$ proposed in Hamilton's Question~4.8.
\end{abstract}

\maketitle

\section{Introduction}\label{sec:intro}

The analytic classification of unibranch curve singularities leads to the classification of complete $\CC$-subalgebras of the power series ring $\CC\ps$. The semigroup of orders of the elements of such a subalgebra is a discrete invariant, and Zariski's moduli problem \cite{Zar86} asks for a description of the singularities with a given semigroup up to analytic equivalence. A finer and more rigid object is the set
\[
\Rs=\{\text{complete $\CC$-subalgebras of $\CC\ps$ with semigroup $\Gamma$}\}
\]
of \emph{embedded} subalgebras. Ishii \cite{Ish80} showed that the corresponding moduli functor is represented by an affine scheme; see the discussion in the introduction of \cite{Ham20}. Hamilton \cite{Ham20} gave an explicit algorithmic description: every element of $\Rs$ has a unique generating set in normal form, and $\Rs$ is the zero set of explicit polynomials in the coefficients of these normal forms \cite[Theorem~2.22]{Ham20}. Hamilton also explained that $\Rs$ parametrizes the ways of gluing a singular point with semigroup $\Gamma$ onto a fixed pointed smooth curve, and that the Zariski moduli space is the quotient of $\Rs$ by $\Aut\CC\ps$ \cite[Sections~1 and~4.3]{Ham20}. In this paper, as in the two questions discussed below, no such quotient is taken: two subalgebras that differ by a change of the parameter $t$ are different points of $\Rs$. This also distinguishes $\Rs$ from crimping spaces, which parametrize singularities of a fixed analytic type; see \cite[Section~1.6]{Bat24} and the references there.

Let $\Gamma=\langle v_0,\dots,v_g\rangle$ with minimal generators $v_0<\dots<v_g$. The variables written $x_0,\dots,x_g$ in \cite{Ham20} are denoted here by $\xi_0,\dots,\xi_g$, and normal-form generators by $\eta_0(t),\dots,\eta_g(t)$, because $x_n$ will denote a coefficient of a generator. An element of $\Rs$ is generated by $g+1$ normal-form elements, but it may be generated by fewer elements \cite[Remark~2.5]{Ham20}. Hamilton denotes by $\Rd{n}$ the subset of subalgebras which can be generated by $n$ elements \cite[Definition~4.4]{Ham20}; this gives a filtration of $\Rs$ by embedding dimension. For semigroups with three generators the plane locus is described in \cite[Section~4.2]{Ham20}, and for arbitrary semigroups the following two questions are posed. In the first one, $\red_{\langle v_0,\dots,v_{i-1}\rangle}$ is the reduction of \cite[Algorithm~2.15]{Ham20}, which removes from an element of $\CC[t]/(t^c)$ the powers lying in $\langle v_0,\dots,v_{i-1}\rangle$ by means of the first $i$ normal-form generators, and $\varphi_R\colon\CC[\xi_0,\dots,\xi_g]\to\CC[t]/(t^c)$ sends $\xi_i$ to the $i$-th normal-form generator $\eta_i(t)$.

\begin{question}[{\cite[Question~4.7]{Ham20}}]\label{q:47}
Can we explicitly write down polynomials $f_d,\dots,f_g\in\CC[\xi_0,\dots,\xi_g]$ such that $R=\CC[\eta_0(t),\dots,\eta_g(t)]$ lies in $\Rd{d}$ if and only if $\ord\red_{\langle v_0,\dots,v_{i-1}\rangle}(\varphi_R(f_i))=v_i$ for all $i\in\{d,\dots,g\}$?
\end{question}

If so, the leading coefficients of these reduced elements would exhibit $\Rd{d}$ as a principal open subvariety of $\Rs$, as Hamilton notes after the question.

\begin{question}[{\cite[Question~4.8]{Ham20}}]\label{q:48}
If $\Rs$ is in bijection with an affine space $\CC^N$, can we always identify $\Rd{d}$ with the space $(\CC^*)^{g+1-d}\times\CC^{N-g+d-1}$?
\end{question}

Here ``in bijection'' and ``identify'' refer to the algebraic structures that are used throughout \cite{Ham20}: the bijections of $\Rs$ with affine spaces in \cite[Section~3]{Ham20} are given by polynomial maps with polynomial inverses, and the plane locus in \cite[Example~4.6]{Ham20} is identified with $\CC^*\times\CC^8$ by such a map. The nonempty loci considered below and their proposed products all have the cardinality of the continuum, so a set-theoretic bijection alone would carry no geometric information.

We study the semigroup
\[
\Gamma=\langle 10,11,34,46,58\rangle,\qquad g=4,
\]
whose conductor is $60$. Write $\Ed{d}\subseteq\Rs$ for the set of subalgebras with embedding dimension exactly $d$, so that $\Rd{d}=\bigcup_{k\le d}\Ed{k}$.

\begin{theorem}\label{thm:main}
Let $\Gamma=\langle 10,11,34,46,58\rangle$. There is an isomorphism of affine varieties
\[
\Phi=(q,A,B,C)\colon\Rs\xrightarrow{\ \sim\ }\AF^{54}\times\AF^3
\]
such that, for every $R\in\Rs$, the minimum number of generators of $R$ equals
\[
\edim R \;=\; 5-\rank\begin{pmatrix}A(R)&C(R)\\0&B(R)\end{pmatrix}.
\]
The same number is the minimum number of topological generators of the corresponding complete subalgebra of $\CC\ps$. Consequently,
\begin{align*}
&\Rd{2}=\emptyset, && \Rd{3}=\Ed{3}\cong\AF^{54}\times\{AB\ne0\},\\
&\Rd{4}\cong\AF^{54}\times(\AF^3\setminus\{0\}), && \Rd{5}=\Rs\cong\AF^{57}.
\end{align*}
\end{theorem}

The coordinates $q$ consist of $54$ of the normal-form coefficients, and $A$, $B$, $C$ are coefficients occurring when two binomial relations are reduced; see Theorem~\ref{thm:coords}. The exact strata $\Ed{4}$ and $\Ed{5}$ are described in Corollary~\ref{cor:strata}.

\begin{corollary}\label{cor:answers}
Let $\Gamma=\langle 10,11,34,46,58\rangle$, so that $g=4$, and take $d=4$.
\begin{enumerate}
\item There is no polynomial $f_4\in\CC[\xi_0,\dots,\xi_4]$ such that $R\in\Rd{4}$ if and only if $\ord\red_{\langle v_0,\dots,v_3\rangle}(\varphi_R(f_4))=v_4$. More strongly, $\Rd{4}$ is an open subvariety of $\Rs$ which is not of the form $\{h\ne0\}$ for a regular function $h$ on $\Rs$, nor a union of two such sets.
\item The space $\Rs$ is isomorphic to $\CC^{57}$, but $\Rd{4}$ is neither isomorphic as a variety nor homeomorphic in the classical topology to $(\CC^*)^{g+1-d}\times\CC^{N-g+d-1}=\CC^*\times\CC^{56}$.
\end{enumerate}
\end{corollary}

Thus Question~\ref{q:47}, in its formulation with fixed polynomials, and Question~\ref{q:48}, when ``identify'' means an isomorphism of varieties or even a homeomorphism, have negative answers. The other strata behave as predicted by Question~\ref{q:48}: $\Rd{3}\cong(\CC^*)^2\times\AF^{55}$ and $\Rd{5}=\Rs\cong\AF^{57}$; the failure occurs exactly at $d=4$. The stratum $\Rd{4}$ is open in $\Rs$, but it is not principal open, as the remark after Question~\ref{q:47} would require; the obstruction is that its complement has codimension three.

The proofs are explicit. Below the conductor, the five normal-form generators produce $31$ monomials but only four coincidences of weights, and all of them are governed by the two binomials $H^4-XZ$ and $H^2Z-XW$ (Lemma~\ref{lem:monomials}). This yields a complete list of five equations for $\Rs$ (Proposition~\ref{prop:criterion}) and a complete list of linear relations in the cotangent space (Proposition~\ref{prop:cotangent}). A weight grading on the coefficient ring, induced by rescaling $t$, controls which coefficients can occur in which equation (Section~\ref{sec:grading}); it allows us to use large polynomials, such as the coordinate $C$, without expanding them.

We do not attempt a classification of semigroups for which the phenomenon occurs. The methods used to describe $\Rs$, namely normal forms, truncation at the conductor and reduction of binomial relations, are those of \cite{Ham20}; see also \cite{AGM17} for closely related division and basis criteria for subalgebras of $K[\![x]\!]$.

\subsection*{Organization}
Section~\ref{sec:prelim} recalls the setting and relates the number of generators to cotangent spaces, including the truncation at the conductor. Section~\ref{sec:semigroup} collects the arithmetic of $\Gamma$, the monomials below the conductor and the weight grading. Section~\ref{sec:equations} determines the equations of $\Rs$ and proves $\Rs\cong\AF^{57}$. Section~\ref{sec:cotangent} computes the cotangent spaces and constructs the coordinates of Theorem~\ref{thm:main}. Section~\ref{sec:family} gives an explicit three-parameter family, and Section~\ref{sec:questions} proves Corollary~\ref{cor:answers}. Appendix~\ref{app:formulas} lists four of the equations in expanded form.

\section{Preliminaries}\label{sec:prelim}

Throughout we work over $\CC$. For $f\in\CC\ps$ or $f\in\CC[t]/(t^c)$ nonzero, $\ord f$ is the smallest exponent with a nonzero coefficient, and $f$ is \emph{monic} if that coefficient is $1$. We set $\ord 0=+\infty$.

\subsection{Complete subalgebras and normal forms}
Let $\Gamma\subseteq\NN$ be a numerical semigroup with conductor $c$, minimally generated by $v_0<\dots<v_g$. If $\widehat R\subseteq\CC\ps$ is a complete $\CC$-subalgebra whose set of orders is $\Gamma$, then $t^c\CC\ps\subseteq\widehat R$ \cite[Proposition~2.1]{Ham20}. Hence $\widehat R$ is determined by its image $R$ in $J_c=\CC[t]/(t^c)$, and $\widehat R\mapsto R$ is a bijection onto the set of $\CC$-subalgebras $R\subseteq J_c$ with
\[
\Gamma_R:=\{\ord f\mid 0\neq f\in R\}\cup\{n\in\NN\mid n\ge c\}=\Gamma.
\]
Following \cite[Section~2]{Ham20} we work with these truncated subalgebras. A tuple $(\eta_0(t),\dots,\eta_g(t))$ in $J_c$ is in \emph{normal form} if each $\eta_i(t)$ is monic of order $v_i$ and all its other exponents with nonzero coefficient are gaps of $\Gamma$ \cite[Definition~2.2]{Ham20}. Every $R$ with $\Gamma_R=\Gamma$ is generated by a unique tuple in normal form \cite[Proposition~2.4]{Ham20}. Taking the coefficients of the normal forms as coordinates identifies $\Rs$ with a Zariski closed subset of $\AF^{M(\Gamma)}$, where $M(\Gamma)=\sum_i\#\{\text{gaps}>v_i\}$ \cite[Theorem~2.22]{Ham20}. We regard $\Rs$ as an affine variety with its reduced structure.

\begin{remark}\label{rem:fixed-t}
The group $\Aut\CC\ps$ of substitutions $t\mapsto\theta(t)$ with $\ord\theta=1$ acts on $\Rs$, and its orbits are the analytic equivalence classes \cite[Section~4.3]{Ham20}. Since the number of generators is invariant under this action, each stratum $\Rd{d}$ is a union of orbits.
\end{remark}

\subsection{Generators and cotangent spaces}
For $R\in\Rs$ let $\m=\{f\in R\mid f=0\text{ or }\ord f>0\}$. Then $R=\CC\oplus\m$, the ideal $\m$ is nilpotent, and $(R,\m)$ is a local Artinian $\CC$-algebra. We put $\edim R=\dim_\CC\m/\m^2$.

\begin{lemma}\label{lem:nakayama}
Let $(R,\m)$ be a finite-dimensional local $\CC$-algebra with $R=\CC\oplus\m$. The minimum number of elements generating $R$ as a $\CC$-algebra is $\dim_\CC\m/\m^2$. Elements $g_1,\dots,g_k\in\m$ generate $R$ if and only if their classes span $\m/\m^2$.
\end{lemma}

\begin{proof}
Subtracting constants does not change the generated algebra, so we may assume that generators lie in $\m$. If $R=\CC[g_1,\dots,g_k]$ with $g_i\in\m$, then $\m$ is spanned by the nonconstant monomials in the $g_i$, and all of them except the $g_i$ lie in $\m^2$. Conversely, suppose that the classes of $g_1,\dots,g_k\in\m$ span $\m/\m^2$ and let $S=\CC[g_1,\dots,g_k]$ and $L=\spn\{g_1,\dots,g_k\}$, so that $\m=L+\m^2$. For $j\ge1$, the space $\m^j$ is spanned by products $y_1\cdots y_j$ with $y_i\in\m$; writing $y_i=\ell_i+n_i$ with $\ell_i\in L$ and $n_i\in\m^2$, we get $y_1\cdots y_j\in\ell_1\cdots\ell_j+\m^{j+1}\subseteq S+\m^{j+1}$. Hence $\m^j\subseteq S+\m^{j+1}$ for all $j\ge1$, and therefore $\m\subseteq S+\m^{j}$ for all $j$. Since $\m$ is nilpotent, $\m\subseteq S$.
\end{proof}

Accordingly,
\[
\Rd{d}=\{R\in\Rs\mid\edim R\le d\},\qquad \Ed{d}=\{R\in\Rs\mid\edim R=d\}.
\]
In \cite[Definition~4.4]{Ham20}, $\Rd{n}$ is written for $3\le n\le g+1$ and the plane locus $\Rd{2}$ is denoted $\mathcal{R}^{\mathrm{plane}}_\Gamma$; generation by $n$ elements includes generation by fewer elements.

The next lemma justifies the truncation at the conductor for questions about generators. It applies to our semigroup by Lemma~\ref{lem:semigroup}.

\begin{lemma}\label{lem:truncation}
Assume that every minimal generator of $\Gamma$ is smaller than $c$. Let $\widehat R\subseteq\CC\ps$ be a complete subalgebra with set of orders $\Gamma$, let $\pi\colon\CC\ps\to J_c$ be the projection, $R=\pi(\widehat R)$ and $I=t^c\CC\ps$. Let $\widehat\m=\{f\in\widehat R\mid f=0\text{ or }\ord f>0\}$ and $\m=\pi(\widehat\m)$.
\begin{enumerate}
\item $I\subseteq\widehat\m^2$, and $\pi$ induces an isomorphism $\widehat\m/\widehat\m^2\cong\m/\m^2$.
\item For $g_1,\dots,g_k\in\widehat\m$, the closure of $\CC[g_1,\dots,g_k]$ in the $t$-adic topology equals $\widehat R$ if and only if $\CC[\pi(g_1),\dots,\pi(g_k)]=R$.
\end{enumerate}
In particular, the minimum number of topological generators of $\widehat R$ equals $\edim R$.
\end{lemma}

\begin{proof}
Since every minimal generator of $\Gamma$ is smaller than $c$, every $n\ge c$ is a sum of two nonzero elements of $\Gamma$: otherwise $n$ would be a minimal generator. For each nonzero $a\in\Gamma$ choose a monic $f_a\in\widehat R$ of order $a$.

(1) For $c\le n<c+v_0$ write $n=a+b$ with $a,b\in\Gamma\setminus\{0\}$; then $\mu_n=f_af_b\in\widehat\m^2$ is monic of order $n$. Moreover $I\subseteq\widehat R$ by \cite[Proposition~2.1]{Ham20}, hence $I\subseteq\widehat\m$, and $f_{v_0}I=t^{c+v_0}\CC\ps$ lies in $\widehat\m^2$, because $f_{v_0}$ is $t^{v_0}$ times a unit. Given $f\in I$, we can subtract a linear combination of $\mu_c,\dots,\mu_{c+v_0-1}$ so that the result lies in $t^{c+v_0}\CC\ps$. Hence $I\subseteq\widehat\m^2$. The map $\widehat\m/\widehat\m^2\to\m/\m^2$ is surjective, and its kernel is $(\pi^{-1}(\m^2)\cap\widehat\m)/\widehat\m^2=(\widehat\m^2+I)/\widehat\m^2=0$.

(2) Let $S$ be the closure of $\CC[g_1,\dots,g_k]$. The set $\widehat R=\pi^{-1}(R)$ is a union of cosets of the open subgroup $I$, hence closed, so $S\subseteq\widehat R$. Since $I$ is open, every element of $S$ is congruent modulo $I$ to an element of $\CC[g_1,\dots,g_k]$, so $\pi(S)=\CC[\pi(g_1),\dots,\pi(g_k)]$. If $S=\widehat R$, this gives $R=\CC[\pi(g_1),\dots,\pi(g_k)]$. Conversely, assume the latter. Then $\pi(S)=R$, so $S+I=\widehat R$, and $S$ contains a monic element $s_a$ of order $a$ for every $a\in\Gamma$ with $a<c$. As in (1), products give monic elements of $S$ of each order in $[c,c+v_0)$, and multiplication by $s_{v_0}$ gives monic elements $s_n\in S$ of every order $n\ge c$. For $f\in I$, successively subtracting multiples of $s_c,s_{c+1},\dots$ produces partial sums in $S$ which converge to $f$; since $S$ is a closed subalgebra, $f\in S$. Thus $S=\widehat R$.

The final assertion follows from (1), (2) and Lemma~\ref{lem:nakayama}.
\end{proof}

\section{The semigroup, monomials and a grading}\label{sec:semigroup}

From now on $\Gamma=\langle10,11,34,46,58\rangle$ and $J=\CC[t]/(t^{60})$.

\begin{lemma}\label{lem:semigroup}
The elements of $\Gamma$ below $60$ are
\[
0,\ 10,\ 11,\ 20,\ 21,\ 22,\ 30,\dots,34,\ 40,\dots,46,\ 50,\dots,58,
\]
twenty-seven in total. Every integer $n\ge60$ lies in $\Gamma$ and $59\notin\Gamma$, so the conductor is $c=60$. The gaps are
\[
1,\dots,9;\quad 12,\dots,19;\quad 23,\dots,29;\quad 35,\dots,39;\quad 47,48,49;\quad 59,
\]
thirty-three in total. The numbers $10,11,34,46,58$ are the minimal generators of $\Gamma$; all of them are smaller than $c$.
\end{lemma}

\begin{proof}
The elements $10a+11b$ below $60$ are $0$, $10$, $11$, $20$, $21$, $22$, $30,\dots,33$, $40,\dots,44$ and $50,\dots,55$. Adding $34$ to those below $26$ gives $34,44,45,54,55,56$; adding $46$ to those below $14$ gives $46,56,57$; and $58$ is the only remaining element obtained from $58$. Sums involving two of $34,46,58$ exceed $59$. This gives the stated list. The integer $59$ is not in the list. The integers $60,\dots,69$ lie in $\Gamma$, as $60=6\cdot10$, $61=50+11$, $62=40+22$, $63=30+33$, $64=20+44$, $65=10+55$, $66=6\cdot11$, $67=33+34$, $68=34+34$ and $69=11+58$; adding multiples of $10$ shows that all $n\ge60$ lie in $\Gamma$. Finally, $34$, $46$ and $58$ are not sums of two smaller nonzero elements of the list, so the five generators are minimal.
\end{proof}

We write $\mathcal{G}$ for the set of gaps and $\Gamma_{<60}=\Gamma\cap[0,59]$.

\subsection{Normal forms and coordinates}
By Lemma~\ref{lem:semigroup}, a tuple in normal form is
\begin{equation}\label{eq:normal}
\begin{gathered}
X=t^{10}+\sum_{n\in\mathcal G,\,n>10}x_nt^n,\qquad
H=t^{11}+\sum_{n\in\mathcal G,\,n>11}h_nt^n,\\
Z=t^{34}+\sum_{n\in\mathcal G,\,n>34}z_nt^n,\qquad
W=t^{46}+\sum_{n\in\mathcal G,\,n>46}w_nt^n,\\
U=t^{58}+u_{59}t^{59}.
\end{gathered}
\end{equation}
The numbers of tail coefficients are $24,24,9,4,1$, so $M(\Gamma)=62$. Note that there is no coefficient $x_{11}$, because $11\in\Gamma$. Let $P$ be the polynomial ring over $\CC$ in the $62$ coordinates $x_n,h_n,z_n,w_n,u_{59}$, so that $\Spec P=\AF^{62}$. We read \eqref{eq:normal} in $J_P=P[t]/(t^{60})$, and for a point $p\in\AF^{62}$ we write $X_p,\dots,U_p\in J$ for the specialized tuple and
\[
R_p=\CC[X_p,H_p,Z_p,W_p,U_p]\subseteq J.
\]
By the results recalled in Section~\ref{sec:prelim}, $\Rs$ is identified with the set of points $p$ such that $\Gamma_{R_p}=\Gamma$.

For $\alpha\in\NN^5$ put $M^\alpha=X^{\alpha_1}H^{\alpha_2}Z^{\alpha_3}W^{\alpha_4}U^{\alpha_5}$ and $\omega(\alpha)=10\alpha_1+11\alpha_2+34\alpha_3+46\alpha_4+58\alpha_5$, the weighted order. Since each generator has leading term $t^v$, the monomial $M^\alpha$ is monic of order $\omega(\alpha)$ at every point when $\omega(\alpha)<60$, and it vanishes in $J$ when $\omega(\alpha)\ge60$.

\begin{lemma}\label{lem:monomials}
There are exactly $31$ monomials $M^\alpha$ with $\omega(\alpha)<60$, including $1$. Their weighted orders are precisely the $27$ elements of $\Gamma_{<60}$. Each element of $\Gamma_{<60}$ is the weighted order of exactly one of them, except
\[
44:\ H^4,\ XZ;\qquad 54:\ XH^4,\ X^2Z;\qquad 55:\ H^5,\ XHZ;\qquad 56:\ H^2Z,\ XW.
\]
\end{lemma}

\begin{proof}
The monomials in $X$ and $H$ of degree $k$ have weighted orders $10k,\dots,11k$, one each; for $k\le5$ this gives $21$ monomials with orders $0$; $10,11$; $20,21,22$; $30,\dots,33$; $40,\dots,44$; $50,\dots,55$, and for $k\ge6$ the order is at least $60$. Monomials containing $Z$ once are $Z$ times $1,X,H,X^2,XH,H^2$, of orders $34,44,45,54,55,56$. Monomials containing $W$ once are $W,XW,HW$, of orders $46,56,57$. Finally there is $U$. Every other monomial has order at least $60$. This gives $21+6+3+1=31$ monomials, and comparing orders gives the stated coincidences.
\end{proof}

\begin{definition}\label{def:basis}
For $k\in\Gamma_{<60}$ let $b_k$ be the unique monomial of weighted order $k$ in Lemma~\ref{lem:monomials}, except that we choose
\[
b_{44}=XZ,\qquad b_{54}=X^2Z,\qquad b_{55}=XHZ,\qquad b_{56}=XW.
\]
Put
\[
F=H^4-XZ,\qquad G=H^2Z-XW .
\]
\end{definition}

Then $XH^4-X^2Z=XF$ and $H^5-XHZ=HF$. In particular $b_0=1$, $b_{10}=X$, $b_{11}=H$, $b_{34}=Z$, $b_{45}=HZ$, $b_{46}=W$, $b_{57}=HW$ and $b_{58}=U$. Every $b_k$ with $k\in\Gamma_{<60}\setminus\{0,10,11,34,46,58\}$ is a monomial of degree at least two.

\begin{lemma}\label{lem:division}
The $P$-module $J_P$ is free with basis $\{b_k\mid k\in\Gamma_{<60}\}\cup\{t^n\mid n\in\mathcal G\}$. Thus every $f\in J_P$ has a unique decomposition
\[
f=\sum_{k\in\Gamma_{<60}}c_k(f)\,b_k+\rho(f),\qquad c_k(f)\in P,\quad \rho(f)=\sum_{n\in\mathcal G}\rho_n(f)\,t^n,\ \rho_n(f)\in P,
\]
which is computed by eliminating, in increasing order of $k\in\Gamma_{<60}$, the coefficient of $t^k$ with the monic element $b_k$. If all coefficients of $f$ below $t^{m}$ vanish, then $c_k(f)=0$ for $k<m$. The decomposition commutes with specialization at any point $p\in\AF^{62}$, where it is the unique decomposition of $f(p)$ with respect to the basis $\{b_k(p)\}\cup\{t^n\}$ of $J$.
\end{lemma}

\begin{proof}
The $60$ elements are monic of the $60$ distinct orders $0,\dots,59$, so the transition matrix to $1,t,\dots,t^{59}$ is unitriangular over $P$. The elimination uses only ring operations in $P$, since each $b_k$ is monic, and it commutes with evaluation at $p$.
\end{proof}

We call $\rho(f)$ the \emph{remainder} of $f$. The polynomials $c_k(f)$ and $\rho_n(f)$ are regular functions on $\AF^{62}$.

\subsection{A weight grading}\label{sec:grading}
Grade $P$ by assigning to each coordinate the difference between its exponent and the order of its generator:
\[
\wt x_n=n-10,\quad \wt h_n=n-11,\quad \wt z_n=n-34,\quad \wt w_n=n-46,\quad \wt u_{59}=1.
\]
All weights are positive; the smallest ones are $\wt h_{12}=\wt z_{35}=\wt w_{47}=\wt u_{59}=1$ and $\wt x_{12}=2$. This grading corresponds to the rescaling $t\mapsto\lambda t$ followed by renormalization of the leading coefficients, which multiplies each coordinate $y$ by $\lambda^{\wt y}$.

\begin{lemma}\label{lem:homog}
A polynomial is homogeneous of degree $\delta$ if it is a linear combination of monomials of degree $\delta$; the zero polynomial is homogeneous of every degree, and a nonzero polynomial is never homogeneous of negative degree.
\begin{enumerate}
\item For every $\alpha$, the coefficient of $t^n$ in $M^\alpha$ is homogeneous of degree $n-\omega(\alpha)$.
\item Let $f\in J_P$ and $\nu\in\ZZ$ be such that the coefficient of $t^n$ in $f$ is homogeneous of degree $n-\nu$ for every $n$. Then $c_k(f)$ is homogeneous of degree $k-\nu$, and $\rho_n(f)$ is homogeneous of degree $n-\nu$.
\end{enumerate}
\end{lemma}

\begin{proof}
(1) The coefficient of $t^n$ in a generator of order $v$ is $1$ for $n=v$, and otherwise it is $0$ or a coordinate of weight $n-v$. For a product, the coefficient of $t^n$ is $\sum_{i+j=n}a_ib_j$, which is homogeneous of degree $(i-\omega)+(j-\omega')=n-\omega-\omega'$. (2) The property is preserved by the elimination of Lemma~\ref{lem:division}: at step $k$, the coefficient $c_k(f)$ is the current coefficient of $t^k$, of degree $k-\nu$, and subtracting $c_k(f)b_k$ changes the coefficient of $t^n$ by a polynomial of degree $(k-\nu)+(n-k)=n-\nu$ by (1).
\end{proof}

\begin{lemma}\label{lem:linear}
Let $Q\in P$ be homogeneous of degree $\delta\ge1$ and let $y$ be a coordinate.
\begin{enumerate}
\item If $\wt y>\delta$, then $Q$ does not involve $y$.
\item If $\wt y=\delta$, then $Q=\gamma y+Q'$ with $\gamma\in\CC$ and $Q'$ not involving $y$. Moreover $\gamma$ is the value of $Q$ at the point where $y=1$ and all other coordinates vanish.
\end{enumerate}
\end{lemma}

\begin{proof}
Since all weights are positive, a monomial containing $y$ has degree at least $\wt y$, with equality only for $y$ itself. The monomials of $Q'$ have positive degree in the other coordinates and therefore vanish at the indicated point.
\end{proof}

We shall use Lemma~\ref{lem:linear}(2) at points where only one coordinate is nonzero. At such a point four of the five generators are monomials $t^{v}$, and the decompositions of Lemma~\ref{lem:division} can be read off directly.

\section{The equations of the parameter space}\label{sec:equations}

\begin{lemma}\label{lem:span}
For every $p\in\AF^{62}$, the algebra $R_p$ is the linear span of the $31$ monomials of Lemma~\ref{lem:monomials} evaluated at $p$, and
\[
R_p=V_p+\spn\{F_p,\ X_pF_p,\ H_pF_p,\ G_p\},\qquad V_p=\spn\{b_k(p)\mid k\in\Gamma_{<60}\}.
\]
Moreover $\dim V_p=27$, and the orders of the nonzero elements of $V_p$ are exactly the elements of $\Gamma_{<60}$.
\end{lemma}

\begin{proof}
$R_p$ is spanned by all monomials in the generators, and those of weighted order at least $60$ vanish in $J$. The four monomials not chosen in Definition~\ref{def:basis} are $H^4=b_{44}+F$, $XH^4=b_{54}+XF$, $H^5=b_{55}+HF$ and $H^2Z=b_{56}+G$. The elements $b_k(p)$ are monic of distinct orders $k$, so they are linearly independent, and a nonzero linear combination has the order of its first nonzero term.
\end{proof}

\begin{proposition}\label{prop:criterion}
A point $p\in\AF^{62}$ lies in $\Rs$ if and only if $\rho(F)(p)=0$ and $\rho(G)(p)=0$. In this case $R_p=V_p$.
\end{proposition}

\begin{proof}
By Lemma~\ref{lem:span}, $V_p\subseteq R_p$ already realizes every order in $\Gamma_{<60}$, so $p\in\Rs$ if and only if no nonzero element of $R_p$ has order in $\mathcal G$.

Assume $p\in\Rs$. The element $\rho(F)(p)=F_p-\sum_kc_k(F)(p)\,b_k(p)$ lies in $R_p$ and its nonzero coefficients sit at gaps; if it were nonzero, its order would be a gap. Hence $\rho(F)(p)=0$, and likewise $\rho(G)(p)=0$.

Conversely, assume that both remainders vanish at $p$; we suppress $p$ from the notation. Then $F,G\in V$. The coefficients of $t^n$ in $F$ vanish for $n\le44$, because $H^4$ and $XZ$ are both monic of order $44$; hence $c_k(F)=0$ for $k\le44$ by Lemma~\ref{lem:division}. Since $\Gamma\cap[45,59]=\{45,46,50,\dots,58\}$,
\[
F=c_{45}(F)\,HZ+c_{46}(F)\,W+\sum_{k=50}^{58}c_k(F)\,b_k .
\]
Multiplying by $X$ and by $H$, the products $Xb_k$ and $Hb_k$ with $k\ge50$ have weighted order at least $60$ and vanish. Therefore
\[
XF=c_{45}(F)\,b_{55}+c_{46}(F)\,b_{56},\qquad HF=c_{45}(F)\,(b_{56}+G)+c_{46}(F)\,b_{57},
\]
using $XHZ=b_{55}$, $XW=b_{56}$, $H^2Z=b_{56}+G$ and $HW=b_{57}$. Both lie in $V$. By Lemma~\ref{lem:span}, $R=V$, and the orders of its nonzero elements are exactly $\Gamma_{<60}$.
\end{proof}

\begin{remark}
Proposition~\ref{prop:criterion} is an instance of the criterion of \cite[Proposition~2.17 and Theorem~2.22]{Ham20}, with a different choice of the monomials used for reduction, and of the binomial criterion for bases in \cite[Theorem~2.9]{AGM17}. We have included the short proof because the identification $R_p=V_p$ and the explicit form of the relations are used below. Since both descriptions characterize the same subset of $\AF^{62}$, they define the same reduced variety $\Rs$.
\end{remark}

\subsection{The low-order terms}
We now make the remainders explicit. Write $f_n$ for the coefficient of $t^n$ in $f\in J_P$, and introduce the names
\[
\delta=c_{45}(F),\quad A=c_{46}(F),\quad C=c_{58}(F),\quad \varepsilon=c_{57}(G),\quad B=c_{58}(G).
\]

\begin{proposition}\label{prop:lowterms}
The remainders are
\[
\rho(F)=r_{47}t^{47}+r_{48}t^{48}+r_{49}t^{49}+r_{59}t^{59},\qquad \rho(G)=q_{59}t^{59},
\]
where
\begin{align}
\delta&=4h_{12}-z_{35},\label{eq:delta}\\
A&=4h_{13}+6h_{12}^2-x_{12}-z_{36}-\delta\,(h_{12}+z_{35})\notag\\
&=4h_{13}+2h_{12}^2-3h_{12}z_{35}+z_{35}^2-x_{12}-z_{36},\label{eq:A}\\
\varepsilon&=2h_{12}+z_{35}-w_{47},\label{eq:eps}\\
B&=2h_{13}+h_{12}^2+2h_{12}z_{35}+z_{36}-x_{12}-w_{48}-\varepsilon\,(h_{12}+w_{47})\notag\\
&=2h_{13}-h_{12}^2+h_{12}z_{35}-h_{12}w_{47}-z_{35}w_{47}+w_{47}^2+z_{36}-x_{12}-w_{48},\label{eq:B}
\end{align}
and
\begin{align}
r_n&=F_n-\delta\,(HZ)_n-A\,w_n\qquad(n=47,48,49),\label{eq:rn}\\
q_{59}&=G_{59}-\varepsilon\,(HW)_{59}-B\,u_{59}.\label{eq:q59}
\end{align}
The polynomials $c_k(F)$ and $\rho_n(F)$ are homogeneous of degrees $k-44$ and $n-44$, and $c_k(G)$, $\rho_n(G)$ are homogeneous of degrees $k-56$ and $n-56$. In particular
\begin{gather*}
\deg A=\deg B=2,\qquad \deg C=14,\qquad \deg q_{59}=3,\\
\deg r_{47}=3,\qquad \deg r_{48}=4,\qquad \deg r_{49}=5,\qquad \deg r_{59}=15.
\end{gather*}
\end{proposition}

\begin{proof}
By Lemma~\ref{lem:homog}(1) the coefficients of $F$ and $G$ are homogeneous of degree $n-44$ and $n-56$, so the degree statements follow from Lemma~\ref{lem:homog}(2). We have $F_n=0$ for $n\le44$ and $G_n=0$ for $n\le56$, as both terms of each binomial are monic of the same order. The gaps above $44$ are $47,48,49,59$ and the only gap above $56$ is $59$, which gives the shape of the remainders. Writing
\begin{gather*}
X=t^{10}(1+x_{12}t^2+\cdots),\qquad H=t^{11}(1+h_{12}t+h_{13}t^2+\cdots),\\
Z=t^{34}(1+z_{35}t+z_{36}t^2+\cdots),
\end{gather*}
we get
\begin{align*}
F&=t^{44}\bigl((4h_{12}-z_{35})t+(4h_{13}+6h_{12}^2-x_{12}-z_{36})t^2+\cdots\bigr),\\
HZ&=t^{45}\bigl(1+(h_{12}+z_{35})t+\cdots\bigr).
\end{align*}
The elimination of Lemma~\ref{lem:division} first subtracts $\delta\,HZ$ with $\delta=F_{45}$, then $A\,W$ with $A$ the resulting coefficient of $t^{46}$; this is \eqref{eq:delta} and \eqref{eq:A}. The next element of $\Gamma$ is $50$, so the coefficients of $t^{47},t^{48},t^{49}$ in $F-\delta HZ-AW$ are never changed again; since $W_n=w_n$ for $n=47,48,49$, this is \eqref{eq:rn}. Similarly
\begin{align*}
H^2Z&=t^{56}\bigl(1+(2h_{12}+z_{35})t+(2h_{13}+h_{12}^2+2h_{12}z_{35}+z_{36})t^2+\cdots\bigr),\\
XW&=t^{56}\bigl(1+w_{47}t+(x_{12}+w_{48})t^2+\cdots\bigr),\\
HW&=t^{57}\bigl(1+(h_{12}+w_{47})t+\cdots\bigr),
\end{align*}
and the elimination subtracts $\varepsilon\,HW$ and $B\,U$, giving \eqref{eq:eps}, \eqref{eq:B} and \eqref{eq:q59}, because $U_{59}=u_{59}$.
\end{proof}

The expansions of $r_{47}$, $r_{48}$, $r_{49}$ and $q_{59}$, with $17$, $22$, $26$ and $24$ terms, are listed in Appendix~\ref{app:formulas}. The polynomials $r_{59}$ and $C$ are large, and we shall use only the structural information of Proposition~\ref{prop:lowterms} and Lemma~\ref{lem:linear} about them.

\subsection{Solving the equations}
Call the $57$ coordinates other than
\[
y_1=z_{37},\quad y_2=z_{38},\quad y_3=w_{49},\quad y_4=z_{39},\quad y_5=z_{49}
\]
the \emph{free coordinates}. They are the $48$ coefficients of $X$ and $H$ together with
\[
z_{35},\ z_{36},\ z_{47},\ z_{48},\ z_{59},\ w_{47},\ w_{48},\ w_{59},\ u_{59}.
\]
Order the five equations as
\[
e_1=r_{47},\quad e_2=r_{48},\quad e_3=q_{59},\quad e_4=r_{49},\quad e_5=r_{59}.
\]

\begin{theorem}\label{thm:affine}
For $k=1,\dots,5$ we have $\deg e_k=\wt y_k$, and
\[
e_k=-y_k+Q_k,\qquad Q_k\in\CC[\text{free coordinates},\,y_1,\dots,y_{k-1}].
\]
Consequently there are homogeneous polynomials $P_k$ in the free coordinates, of degree $\wt y_k$, such that
\[
\Rs=\{p\in\AF^{62}\mid y_k(p)=P_k(p)\ \text{for }k=1,\dots,5\},
\]
and the projection to the free coordinates is an isomorphism $\Rs\cong\AF^{57}$.
\end{theorem}

\begin{proof}
The weights are $\wt z_{37}=3$, $\wt z_{38}=4$, $\wt w_{49}=3$, $\wt z_{39}=5$, $\wt z_{49}=15$, which equal the degrees $3,4,3,5,15$ of Proposition~\ref{prop:lowterms}. By Lemma~\ref{lem:linear}(2), $e_k=\gamma_ky_k+Q_k'$ with $Q_k'$ not involving $y_k$, and $\gamma_k$ is obtained by evaluating at the point where $y_k=1$ and all other coordinates vanish.

At the point $z_n=1$ with $n\in\{37,38,39,49\}$ we have $X=t^{10}$, $H=t^{11}$, $Z=t^{34}+t^n$, so $F=-t^{n+10}$, which is a gap monomial for $n+10\in\{47,48,49,59\}$. Hence $\rho(F)=-t^{n+10}$ and the coefficient of $y_k$ in $r_{n+10}$ is $-1$ for $k=1,2,4,5$. At the point $w_{49}=1$ we have $W=t^{46}+t^{49}$ and all other generators monomial, so $G=t^{56}-(t^{56}+t^{59})=-t^{59}$, and the coefficient of $y_3$ in $q_{59}$ is $-1$.

It remains to show that $e_k$ does not involve $y_{k+1},\dots,y_5$. Since $\wt z_{38}=4>\deg e_1$, $\wt z_{39}=5>\deg e_1,\deg e_2,\deg e_3$, and $\wt z_{49}=15>\deg e_1,\dots,\deg e_4$, Lemma~\ref{lem:linear}(1) leaves only the occurrence of $w_{49}$ in $e_1=r_{47}$ and $e_2=r_{48}$. In \eqref{eq:rn}, the polynomials $F_n$, $(HZ)_n$, $\delta$ and $A$ involve no coefficient of $W$, so the only coefficient of $W$ in $r_n$ is $w_n$. Hence $w_{49}$ occurs in neither $r_{47}$ nor $r_{48}$.

By Proposition~\ref{prop:criterion}, $\Rs$ is the zero set of $e_1,\dots,e_5$. The system $y_k=Q_k$ is triangular: substituting recursively gives polynomials $P_k$ in the free coordinates, homogeneous of degree $\wt y_k$ because the $Q_k$ are, and $\Rs$ is the graph of the polynomial map $(P_1,\dots,P_5)$. The projection to the free coordinates is therefore an isomorphism, with inverse given by the $P_k$.
\end{proof}

\begin{remark}\label{rem:scheme}
Put $I_k=(e_1,\dots,e_k)$ and $I'_k=(y_1-P_1,\dots,y_k-P_k)$ in $P$. If $I_{k-1}=I'_{k-1}$, then $Q_k\equiv P_k$ modulo $I'_{k-1}$, because $P_k$ is obtained from $Q_k$ by substituting $P_j$ for $y_j$, $j<k$; hence $e_k\equiv-(y_k-P_k)$ and $I_k=I'_k$. By induction $I_5=I'_5$, and $P/I'_5$ is the polynomial ring in the $57$ free coordinates. Thus the closed subscheme of $\AF^{62}$ defined by $e_1,\dots,e_5$ is reduced and isomorphic to $\AF^{57}$.
\end{remark}

\section{Cotangent spaces and global coordinates}\label{sec:cotangent}

\begin{proposition}\label{prop:cotangent}
Let $p\in\Rs$ and $R=R_p$. In $\m/\m^2$ the classes of $X,H,Z,W,U$ span, and the space of linear relations among them is spanned by
\[
A\,[W]+C\,[U]=0\qquad\text{and}\qquad B\,[U]=0,
\]
with $A,B,C$ evaluated at $p$. Consequently
\[
\edim R=5-\rank\begin{pmatrix}A&C\\0&B\end{pmatrix}.
\]
\end{proposition}

\begin{proof}
Let $\Gamma_2=\Gamma_{<60}\setminus\{0,10,11,34,46,58\}$, a set of $21$ elements, and $N=\spn\{b_k\mid k\in\Gamma_2\}$. By Proposition~\ref{prop:criterion}, $R=V$ has basis $\{b_k\}$, so
\[
\m=N\oplus\spn\{X,H,Z,W,U\}.
\]
The ideal $\m^2$ is spanned by the products of two nonconstant monomials, that is, by the $25$ monomials of degree at least two among the $31$ monomials of Lemma~\ref{lem:monomials}. These are the $21$ elements $b_k$, $k\in\Gamma_2$, together with $H^4=b_{44}+F$, $XH^4=b_{54}+XF$, $H^5=b_{55}+HF$ and $H^2Z=b_{56}+G$. Hence
\[
\m^2=N+\spn\{F,\ XF,\ HF,\ G\}.
\]
We compute these four elements modulo $N$, using the decompositions in the proof of Proposition~\ref{prop:criterion}. Since $45,50,\dots,57\in\Gamma_2$,
\begin{align*}
F&\equiv A\,W+C\,U, & G&=\varepsilon\,b_{57}+B\,U\equiv B\,U,\\
XF&\equiv0, & HF&\equiv\delta\,G\equiv\delta B\,U \qquad\pmod N.
\end{align*}
Therefore $\m^2=N\oplus\spn\{AW+CU,\ BU\}$, and
\[
\m/\m^2\cong\spn\{X,H,Z,W,U\}\big/\spn\{AW+CU,\ BU\}.
\]
The dimension of the denominator is the rank of $\left(\begin{smallmatrix}A&C\\0&B\end{smallmatrix}\right)$.
\end{proof}

In particular $\edim R\ge3$ for every $R\in\Rs$, so $\Rd{2}=\emptyset$: by Lemma~\ref{lem:truncation}, none of the corresponding complete subalgebras of $\CC\ps$ is the local ring of a plane branch. This is consistent with the criterion recalled in \cite[Section~4.2]{Ham20}, which fails for $\Gamma$ already because $\gcd(10,11)=1$.

\begin{lemma}\label{lem:C}
On $\Rs$, expressed in the free coordinates,
\[
A=a-z_{36},\qquad B=b-w_{48},\qquad C=-z_{48}+\Psi,
\]
where $a$ is a polynomial in $x_{12}$, $h_{12}$, $h_{13}$, $z_{35}$; $b$ is a polynomial in $x_{12}$, $h_{12}$, $h_{13}$, $z_{35}$, $z_{36}$, $w_{47}$; and $\Psi$ is a polynomial in the free coordinates other than $z_{48}$.
\end{lemma}

\begin{proof}
The statements about $A$ and $B$ are \eqref{eq:A} and \eqref{eq:B}; these formulas involve no eliminated coordinate. The polynomial $C=c_{58}(F)\in P$ is homogeneous of degree $14=\wt z_{48}$. By Lemma~\ref{lem:linear}, it does not involve $z_{49}$, of weight $15$, and $C=\gamma z_{48}+C'$ with $C'$ free of $z_{48}$. At the point where $z_{48}=1$ and all other coordinates vanish, $F=-t^{58}=-U$, so $\gamma=-1$. On $\Rs$ we substitute $y_k=P_k$ for $k\le4$; these $P_k$ have degrees $3,4,3,5<14$, so they do not involve $z_{48}$ by Lemma~\ref{lem:linear}(1). Hence $\Psi:=C'|_{y_k=P_k}$ does not involve $z_{48}$.
\end{proof}

\begin{theorem}\label{thm:coords}
Let $q$ denote the $54$ free coordinates other than $z_{36},w_{48},z_{48}$, namely the $48$ coefficients of $X$ and $H$ together with $z_{35},z_{47},z_{59},w_{47},w_{59},u_{59}$. Then
\[
\Phi=(q,A,B,C)\colon\Rs\longrightarrow\AF^{54}\times\AF^3
\]
is an isomorphism of affine varieties, and $\edim R=5-\rank\left(\begin{smallmatrix}A&C\\0&B\end{smallmatrix}\right)$ for every $R\in\Rs$.
\end{theorem}

\begin{proof}
By Theorem~\ref{thm:affine} it suffices to show that $(q,z_{36},w_{48},z_{48})\mapsto(q,A,B,C)$ is an automorphism of $\AF^{57}$. By Lemma~\ref{lem:C} it is given by
\[
A=a(q)-z_{36},\qquad B=b(q,z_{36})-w_{48},\qquad C=-z_{48}+\Psi(q,z_{36},w_{48}),
\]
and its inverse is
\[
z_{36}=a(q)-A,\qquad w_{48}=b\bigl(q,a(q)-A\bigr)-B,\qquad z_{48}=\Psi(q,z_{36},w_{48})-C,
\]
with $z_{36}$ and $w_{48}$ substituted in the last formula. Both maps are polynomial. The formula for $\edim$ is Proposition~\ref{prop:cotangent}.
\end{proof}

\begin{corollary}\label{cor:strata}
Under $\Phi$, the embedding dimension strata are as follows.
\begin{center}
\begin{tabular}{@{}lll@{}}
\toprule
$d$ & exactly $d$: $\Ed{d}$ & at most $d$: $\Rd{d}$\\
\midrule
$\le2$ & $\emptyset$ & $\emptyset$\\
$3$ & $\AF^{54}\times\{AB\neq0\}$ & $\AF^{54}\times\{AB\neq0\}$\\
$4$ & $\AF^{54}\times\bigl(\{AB=0\}\setminus\{0\}\bigr)$ & $\AF^{54}\times(\AF^3\setminus\{0\})$\\
$5$ & $\AF^{54}\times\{0\}$ & $\AF^{54}\times\AF^3$\\
\bottomrule
\end{tabular}
\end{center}
In particular $\Rd{3}\cong(\CC^*)^2\times\AF^{55}$ is dense in $\Rs$, $\Ed{5}\cong\AF^{54}$ is closed of codimension three, and $\Rd{4}$ is a smooth irreducible open subset of $\Rs$. The stratum $\Ed{4}$ has two irreducible components, $\AF^{54}\times(\{A=0\}\setminus\{0\})$ and $\AF^{54}\times(\{B=0\}\setminus\{0\})$, each of dimension $56$; they meet along $\AF^{54}\times\{A=B=0,\ C\ne0\}$, of dimension $55$, which is the singular locus of $\Ed{4}$.
\end{corollary}

\begin{proof}
The matrix $\left(\begin{smallmatrix}A&C\\0&B\end{smallmatrix}\right)$ has determinant $AB$; it has rank $2$ if $AB\ne0$, rank $0$ if $A=B=C=0$, and rank $1$ otherwise. The remaining assertions concern the subsets $\{AB\ne0\}$, $\{AB=0\}\setminus\{0\}$ and $\{0\}$ of $\AF^3$, the second of which is the union of two coordinate planes without the origin.
\end{proof}

Theorem~\ref{thm:main} follows from Theorems~\ref{thm:affine} and~\ref{thm:coords}, Corollary~\ref{cor:strata} and Lemmas~\ref{lem:nakayama} and~\ref{lem:truncation}.

\section{An explicit three-parameter family}\label{sec:family}

The following family is a closed subvariety $\AF^3\subseteq\Rs$ which meets each of the nonempty strata $\Ed{3}$, $\Ed{4}$, $\Ed{5}$. It also gives an independent check of the formulas for $A$, $B$, $C$, through identities which hold exactly in $\CC\ps$.

\begin{example}\label{ex:family}
For $(s_1,s_2,s_3)\in\CC^3$ put
\begin{gather*}
X=t^{10},\qquad H=t^{11},\qquad Z=t^{34}+s_1t^{36}+s_1s_2\,t^{38}+s_3\,t^{48},\\
W=t^{46}+s_2\,t^{48},\qquad U=t^{58}.
\end{gather*}
This tuple is in normal form, since $36,38,48$ are gaps larger than $34$ and $48$ is a gap larger than $46$. In $\CC\ps$ one has
\[
F=H^4-XZ=-s_1W-s_3U,\qquad G=H^2Z-XW=(s_1-s_2)\,U+s_1s_2\,X^6+s_3X^7.
\]
Modulo $t^{60}$ the last two terms vanish, so $\rho(F)=\rho(G)=0$ and the point, which we denote by $\kappa(s_1,s_2,s_3)$, lies in $\Rs$ by Proposition~\ref{prop:criterion}. By uniqueness of the decomposition, $c_{46}(F)=-s_1$, $c_{58}(F)=-s_3$, $c_{57}(G)=0$ and $c_{58}(G)=s_1-s_2$, that is,
\[
(A,B,C)\bigl(\kappa(s_1,s_2,s_3)\bigr)=(-s_1,\ s_1-s_2,\ -s_3),
\]
in agreement with \eqref{eq:A} and \eqref{eq:B}. All coordinates of $\kappa(s_1,s_2,s_3)$ in $q$ vanish, so $\Phi\circ\kappa(s_1,s_2,s_3)=(0,-s_1,s_1-s_2,-s_3)$ and $\kappa$ is an isomorphism of $\AF^3$ onto the closed subvariety $\Phi^{-1}(\{0\}\times\AF^3)$. By Theorem~\ref{thm:coords},
\[
\edim\kappa(s_1,s_2,s_3)=
\begin{cases}
3,& s_1(s_1-s_2)\ne0,\\
5,& s_1=s_2=s_3=0,\\
4,&\text{otherwise.}
\end{cases}
\]
For example, $\kappa(1,0,0)$ is generated by $X,H,Z$, since $W=XZ-H^4$ and $U=H^2Z-XW$; the point $\kappa(1,1,0)$ has embedding dimension four, with $W=XZ-H^4$; and $\kappa(0,0,0)$ is the monomial algebra.
\end{example}

At $\kappa(1,1,0)$ the orders $10,11,34,46,58$ of the normal-form generators minimally generate $\Gamma$, but the generators themselves do not minimally generate the algebra. This is an instance of the phenomenon noted in \cite[Remark~2.5]{Ham20}.

\section{Hamilton's questions}\label{sec:questions}

\subsection{Question~\ref{q:47}}
Let $T=\langle10,11,34,46\rangle$. By Lemma~\ref{lem:semigroup}, $T\cap[0,57]=\Gamma\cap[0,57]$ and $58\notin T$. For $f\in\CC[\xi_0,\dots,\xi_4]$ and $p\in\Rs$, let $\varphi_p(f)=f(X_p,H_p,Z_p,W_p,U_p)\in J$. The reduction $\red_{\langle v_0,\dots,v_3\rangle}$ of \cite[Algorithm~2.15]{Ham20} proceeds as follows: for each $n\in T$ with $0<n<60$ a monic monomial $m_n$ in $\xi_0,\dots,\xi_3$ of weighted degree $n$ is fixed, and, in increasing order of $n$, the coefficient of $t^n$ is removed by subtracting the appropriate multiple of $m_n(X_p,H_p,Z_p,W_p)$. The argument below uses only these features and not the particular choice of $m_n$.

\begin{lemma}\label{lem:reduction}
Let $f\in\CC[\xi_0,\dots,\xi_4]$ with $f(0)=0$, and put $\sigma_p=\red_{\langle v_0,\dots,v_3\rangle}(\varphi_p(f))$ for $p\in\Rs$. There is a polynomial $h_f\in P$ such that, for every $p\in\Rs$,
\[
\ord\sigma_p=58\iff h_f(p)\ne0 .
\]
If $f(0)\neq0$, then $\ord\sigma_p=0$ for every $p\in\Rs$.
\end{lemma}

\begin{proof}
Each $m_n(X,H,Z,W)$ is monic of order $n$ with coefficients in $P$, so every step of the reduction is a polynomial operation, and the coefficient $h_f$ of $t^{58}$ in the output is a polynomial in the coordinates. The constant term of $\varphi_p(f)$ is $f(0)$, and the reduction only subtracts elements of positive order; this proves the last assertion. Assume $f(0)=0$. Then $\sigma_p\in R_p$ has no constant term, and its coefficients at all $n\in T\cap[1,59]$ vanish, since later steps only change higher coefficients. As $p\in\Rs$, a nonzero $\sigma_p$ has order in $\Gamma\cap[1,59]$. The elements of $\Gamma\cap[1,57]$ lie in $T$, and $59$ is a gap; hence either $\sigma_p=0$ or $\ord\sigma_p=58$, and the latter holds exactly when $h_f(p)\ne0$.
\end{proof}

\begin{lemma}\label{lem:twovar}
No polynomial $\phi\in\CC[s_1,s_2]$ vanishes at exactly one point of $\CC^2$.
\end{lemma}

\begin{proof}
Suppose $\phi(a,b)=0$, so $\phi$ is not a nonzero constant. If $\phi=0$ there is nothing to prove. If $\phi$ involves $s_2$, write $\phi=\sum_{j=0}^m \phi_j(s_1)s_2^j$ with $m\ge1$ and $\phi_m\ne0$; for each of the infinitely many $s_1$ with $\phi_m(s_1)\ne0$ there is a root $s_2$. If $\phi$ does not involve $s_2$, then $\phi(a,s_2)=0$ for all $s_2$.
\end{proof}

\begin{theorem}\label{thm:q47}
Let $\Gamma=\langle 10,11,34,46,58\rangle$.
\begin{enumerate}
\item There is no polynomial $f\in\CC[\xi_0,\dots,\xi_4]$ such that, for all $p\in\Rs$,
\[
p\in\Rd{4}\iff\ord\red_{\langle v_0,\dots,v_3\rangle}(\varphi_p(f))=58 .
\]
\item There is no regular function $h$ on $\Rs$ with $\Rd{4}=\{h\neq0\}$.
\item There are no regular functions $h_1,h_2$ on $\Rs$ such that
\[
\Rd{4}=\{h_1\neq0\}\cup\{h_2\ne0\}.
\]
\end{enumerate}
On the other hand, $\Rd{4}=\{A\ne0\}\cup\{B\ne0\}\cup\{C\ne0\}$.
\end{theorem}

\begin{proof}
(2) Let $h$ be a regular function on $\Rs$ and consider the plane $\lambda(s_1,s_2)=\kappa(s_1,s_2,0)$ of Example~\ref{ex:family}. Then $h\circ\lambda\in\CC[s_1,s_2]$, and by Example~\ref{ex:family}, $\lambda(s_1,s_2)\in\Rd{4}$ if and only if $(s_1,s_2)\ne(0,0)$. If $\Rd{4}=\{h\ne0\}$, then $h\circ\lambda$ vanishes exactly at the origin, contradicting Lemma~\ref{lem:twovar}.

(1) If $f(0)\ne0$, the condition never holds by Lemma~\ref{lem:reduction}, while $\Rd{4}\ne\emptyset$. If $f(0)=0$, Lemma~\ref{lem:reduction} gives $\{p\mid\ord\sigma_p=58\}=\{h_f\ne0\}$, which differs from $\Rd{4}$ by (2).

(3) Under $\Phi$, the complement of $\Rd{4}$ is $\Ed{5}=\AF^{54}\times\{0\}$, which is irreducible and nonempty of codimension three in $\AF^{57}$. The complement of $\{h_1\ne0\}\cup\{h_2\ne0\}$ is the zero set of $h_1,h_2$. If it is nonempty, each of its irreducible components corresponds to a minimal prime over $(h_1,h_2)$ in the polynomial ring $\CC[\AF^{57}]$, which has height at most two by Krull's height theorem \cite[Theorem~13.5]{Mat86}; hence each component has codimension at most two. This is impossible.

Finally, the description of $\Rd{4}$ by $A$, $B$, $C$ is Corollary~\ref{cor:strata}.
\end{proof}

Theorem~\ref{thm:q47} proves Corollary~\ref{cor:answers}(1). It concerns the formulation of Question~\ref{q:47} with polynomials $f_d,\dots,f_g$ fixed independently of $R$. By Theorem~\ref{thm:q47}(2) and~(3), at least three principal open subsets are needed to cover $\Rd{4}$, and the coordinates $A$, $B$, $C$ show that three suffice.

\subsection{Question~\ref{q:48}}

\begin{lemma}\label{lem:nonaffine}
Let $m\ge0$, let $\CC[q,A,B,C]$ be the coordinate ring of $\AF^m\times\AF^3$, and let $\mathcal V=\AF^m\times(\AF^3\setminus\{0\})$.
\begin{enumerate}
\item The ring of regular functions on $\mathcal V$ is $\CC[q,A,B,C]$; in particular its units are the nonzero constants.
\item $\mathcal V$ is not affine.
\item In the classical topology, $\mathcal V$ is simply connected.
\end{enumerate}
\end{lemma}

\begin{proof}
(1) Write $S=\CC[q,A,B,C]$. The open sets $\{A\ne0\}$, $\{B\ne0\}$, $\{C\ne0\}$ cover $\mathcal V$ and have rings of regular functions $S_A$, $S_B$, $S_C$, so the ring of regular functions on $\mathcal V$ is $S_A\cap S_B\cap S_C$ inside the fraction field of $S$. If $u=a/A^i=b/B^j$ with $a,b\in S$, then $aB^j=bA^i$; since $S$ is a unique factorization domain and $A,B$ are non-associate primes, $A^i$ divides $a$, so $u\in S$.

(2) If $\mathcal V$ were affine, the inclusion $\mathcal V\hookrightarrow\AF^{m+3}$ would be a morphism of affine varieties inducing an isomorphism of coordinate rings by (1), hence an isomorphism of varieties \cite[Corollary~I.3.7]{Har77}. It is not surjective, a contradiction. (Compare \cite[Exercise~I.3.6]{Har77}.)

(3) $\CC^3\setminus\{0\}$ deformation retracts onto the sphere $S^5$, so $\mathcal V$ is homotopy equivalent to $S^5$, which is simply connected \cite[Proposition~1.14]{Hat02}.
\end{proof}

\begin{theorem}\label{thm:q48}
The space $\Rs$ is isomorphic to $\AF^{57}$, so Question~\ref{q:48} applies with $g=4$ and $N=57$; for $d=4$ it proposes $\CC^*\times\CC^{56}$. However, $\Rd{4}\cong\AF^{54}\times(\AF^3\setminus\{0\})$ is not affine, has only constant units, and is simply connected. Therefore $\Rd{4}$ is not isomorphic to $\CC^*\times\CC^{56}$ as an algebraic variety, and not homeomorphic to it in the classical topology.
\end{theorem}

\begin{proof}
The description of $\Rd{4}$ is Corollary~\ref{cor:strata}, and its properties are Lemma~\ref{lem:nonaffine}. The variety $\CC^*\times\CC^{56}$ is affine, its coordinate ring contains the nonconstant unit given by the coordinate of $\CC^*$, and its fundamental group is $\pi_1(S^1)\cong\ZZ$ \cite[Theorem~1.7]{Hat02}, since it deformation retracts onto a circle.
\end{proof}

Theorem~\ref{thm:q48} proves Corollary~\ref{cor:answers}(2).

\begin{remark}
For the other values of $d$ the prediction of Question~\ref{q:48} is correct for this semigroup. For $d=3$ it gives $(\CC^*)^2\times\CC^{55}$, which is isomorphic to $\Rd{3}$ by Corollary~\ref{cor:strata}, and for $d=5=g+1$ it gives $\CC^{57}=\Rs$. Only $\Rd{4}$, whose complement $\Ed{5}$ has codimension three, fails to have the predicted form.
\end{remark}

\clearpage
\appendix

\section{Expanded equations}\label{app:formulas}

Using \eqref{eq:A}, \eqref{eq:B}, \eqref{eq:rn} and \eqref{eq:q59}, the equations $r_{47},r_{48},r_{49},q_{59}$ of Theorem~\ref{thm:affine} are
\[
r_n=R_n-A\,w_n\quad(n=47,48,49),
\]
where
\begin{align*}
R_{47}&=-z_{37}-x_{13}+4h_{14}+8h_{12}h_{13}+4h_{12}^3-4h_{12}^2z_{35}+h_{12}z_{35}^2-4h_{12}z_{36}\\
&\qquad+h_{13}z_{35}-x_{12}z_{35}+z_{35}z_{36},\\[2pt]
R_{48}&=-z_{38}-x_{14}+4h_{15}+8h_{12}h_{14}+6h_{13}^2+12h_{12}^2h_{13}+h_{12}^4-4h_{12}h_{13}z_{35}\\
&\qquad+h_{13}z_{35}^2+h_{14}z_{35}-4h_{12}^2z_{36}+h_{12}z_{35}z_{36}-4h_{12}z_{37}+z_{35}z_{37}\\
&\qquad-x_{12}z_{36}-x_{13}z_{35},\\[2pt]
R_{49}&=-z_{39}-x_{15}+4h_{16}+8h_{12}h_{15}+12h_{13}h_{14}+12h_{12}^2h_{14}+12h_{12}h_{13}^2+4h_{12}^3h_{13}\\
&\qquad-4h_{12}h_{14}z_{35}+h_{14}z_{35}^2+h_{15}z_{35}-4h_{12}h_{13}z_{36}+h_{13}z_{35}z_{36}-4h_{12}^2z_{37}\\
&\qquad+h_{12}z_{35}z_{37}-4h_{12}z_{38}+z_{35}z_{38}-x_{12}z_{37}-x_{13}z_{36}-x_{14}z_{35},
\end{align*}
and
\begin{align*}
q_{59}&=-w_{49}+z_{37}-x_{13}+2h_{14}+h_{13}z_{35}+h_{13}w_{47}+h_{12}^2z_{35}-2h_{12}^2w_{47}+h_{12}w_{47}^2\\
&\qquad-h_{12}w_{47}z_{35}+2h_{12}z_{36}-2h_{12}w_{48}-w_{48}z_{35}+w_{47}w_{48}-w_{47}x_{12}-B\,u_{59}.
\end{align*}
The pivot terms $-z_{37}$, $-z_{38}$, $-w_{49}$, $-z_{39}$ of Theorem~\ref{thm:affine} are visible, and one checks directly that each of these polynomials is homogeneous of the degree given in Proposition~\ref{prop:lowterms}.

\section*{Use of AI-assisted tools}
AI-assisted tools were used in the research and preparation of this manuscript.

\begin{thebibliography}{AGM17}

\bibitem[AGM17]{AGM17}
A.~Assi, P.~A. Garc\'{\i}a-S\'anchez, and V.~Micale, \emph{Bases of subalgebras of $\mathbb{K}[\![x]\!]$ and $\mathbb{K}[x]$}, J. Symbolic Comput. \textbf{79} (2017), 4--22. arXiv:1412.4089.

\bibitem[Bat24]{Bat24}
L.~Battistella, \emph{Gorenstein curve singularities of genus three}, Ann. Univ. Ferrara Sez. VII Sci. Mat. \textbf{70} (2024), 655--680. \url{https://doi.org/10.1007/s11565-024-00495-8}.

\bibitem[Ham20]{Ham20}
E.~Hamilton, \emph{Classifying complete $\mathbb{C}$-subalgebras of $\mathbb{C}[\![t]\!]$}, Manuscripta Math. \textbf{163} (2020), 437--462. \url{https://doi.org/10.1007/s00229-019-01162-5}.

\bibitem[Har77]{Har77}
R.~Hartshorne, \emph{Algebraic Geometry}, Graduate Texts in Mathematics, vol.~52, Springer-Verlag, New York, 1977.

\bibitem[Hat02]{Hat02}
A.~Hatcher, \emph{Algebraic Topology}, Cambridge University Press, Cambridge, 2002.

\bibitem[Ish80]{Ish80}
S.~Ishii, \emph{Moduli of subrings of a local ring}, J. Algebra \textbf{67} (1980), 504--516.

\bibitem[Mat86]{Mat86}
H.~Matsumura, \emph{Commutative Ring Theory}, Cambridge Studies in Advanced Mathematics, vol.~8, Cambridge University Press, Cambridge, 1986.

\bibitem[Zar86]{Zar86}
O.~Zariski, \emph{Le probl\`eme des modules pour les branches planes}, with an appendix by B.~Teissier, Hermann, Paris, 1986.

\end{thebibliography}

\end{document}
